\documentclass[8pt]{beamer}
\usepackage{tikz}
\usepackage{tikz-cd}
\usepackage{qtree}
\usepackage{amsfonts}
\usepackage{amssymb}
\usepackage{amsmath}
\usepackage{amscd}
\usepackage{amsthm}
% \usepackage{mathpartir}
\usepackage{bbm}
\usepackage[mathcal]{euscript}
% \usepackage[utf8]{inputenc}
\setbeamertemplate{navigation symbols}{}%remove navigation symbols
%\usepackage[T1]{fontenc}
\usepackage{lmodern} \title{The CBH theorem}
\author{Hans Halvorson}
\date[November 2013]{November 12, 2013} \usetheme{Copenhagen}
\begin{document}
\def\2{\mathcal}
\def\7{\mathbb}
\newenvironment{proenv}{\only{\setbeamercolor{local structure}{fg=green}}}{}
\newenvironment{conenv}{\only{\setbeamercolor{local structure}{fg=red}}}{}  
\newenvironment{neuenv}{\only{\setbeamercolor{local
      structure}{fg=yellow}}}{}

%% A theory is _nothing less_ than a category ... but it may be more.

\newcommand{\3}{\mathfrak}

\begin{frame}[plain]
  \titlepage
\end{frame}

\begin{frame}

  \textbf{Goal:} Distinguish between classical and quantum theory in
  terms of their distinctive information-theoretic characteristics.

  \bigskip \textbf{Constraint:} To \textbf{prove} a characterization
  theorem, we must assume a background mathematical framework that is
  general enough to accommodate both classical and quantum sytems.

\end{frame}

\begin{frame}{classical mechanics with $C^*$-algebras}

  If $X$ is the phase space of a classical system, then its
  observables are represented by elements of $C(X)$, the continuous
  functions from $X$ to $\7C$ that are constant at infinity.

  \bigskip $C(X)$ has pointwise sums, products, and complex
  conjugates.  Equipped with the supremum norm, it is an
  \textbf{abelian} or \textbf{commutative} $C^*$-algebra.

  \bigskip The pure states of $C(X)$ correspond to points in $X$.  The
  mixed states of $C(X)$ correspond to probability measures on $X$.

\end{frame}

\begin{frame}{composite classical systems}

If $Y$ is the phase space of another classical system, then:
$$ C(X\times Y) \: \simeq \: C(X)\otimes C(Y) .$$

\bigskip \textbf{Fact:} No states are entangled across $C(X)\otimes
C(Y)$.  The Bell inequality is satisfied in all states.  (Easy proof
using the fact that $C(X)$ and $C(Y)$ are abelian.)

\end{frame}

\begin{frame}{cloning}

  The system $\2A\otimes\2A$ allows \textbf{cloning} just in case
  there is a state $\omega$ of $\2A$ and a CP map
  $T:\2A\otimes\2A\to\2A\otimes\2A$ such that
$$ T^*(\rho\otimes\omega ) \: = \: \rho\otimes\rho ,$$
for all states $\rho$ of $\2A$.

\bigskip \textbf{Fact:} The system $C(X)\otimes C(X)$ allows cloning.

\smallskip \textbf{Proof:} Verify that any homomorphism $f:Y\to Z$
gives rise to a CP map $f^*:C(Z)\to C(Y)$.  Then note that the map
$f(\langle x,z\rangle )=\langle x,x\rangle$ is a homomorphism.



\end{frame}



\begin{frame}{all classical systems allow cloning}

  According to the \textbf{Gelfand representation theorem}, a
  $C^*$-algebra $\2A$ is abelian iff $\2A\simeq C(X)$ for some
  topological space $X$.

  \bigskip \textbf{Corollary:} If $\2A$ is commutative then
  $\2A\otimes\2A$ allows cloning.



\end{frame}

\begin{frame}{no quantum system allows cloning}

  \textbf{New use for an old concept:} If $K$ is a convex set, and
  $x,y\in K$, then the \textbf{affine ratio} between $x,y$, written
  $p(x|y)$, is the largest value of $f(y)$ where $f$  ... 

  \bigskip \textbf{Fact:} The set of states of a classical system is a
  simplex; hence the affine ratio between any two pure states is $0$.

  \bigskip \textbf{Fact:} If $x,y$ are quantum states with affine
  ratio $0$, then there is a continuous path of quantum states from
  $x$ to $y$, achieving all intermediate affine ratios.






\end{frame}


\end{document}